\documentclass[twocolumn,10pt]{article}
\usepackage[margin=1.5cm]{geometry}
\usepackage{amsmath, amssymb}
\usepackage[colorlinks,linkcolor=blue,urlcolor=blue,citecolor=blue]{hyperref}
\usepackage{stfloats}

\title{Mathematics of Spherical and Flat 2D and 3D Worlds}
\author{Tobias Kaiser, 2026}
\date{}

\begin{document}
\maketitle

\begin{table*}[!b]
\centering
\renewcommand{\arraystretch}{1.2}
\begin{tabular}{|l|c|c|c|c|}
\hline
    & \textbf{Flat 2D} & \textbf{Flat 3D} & \textbf{Spherical 2D} & \textbf{Spherical 3D} \\
\hline
Embedding & $\mathbb{R}^2$ & $\mathbb{R}^3$ & $\mathbb{R}^3$ & $\mathbb{R}^4$ \\
\hline
Curvature & \multicolumn{2}{c|}{Zero} & \multicolumn{2}{c|}{Positive} \\
\hline
Extent & Infinite & Infinite & Finite (area $4\pi$) & Finite (vol.\ $2\pi^2$) \\
\hline
Maximum distance & \multicolumn{2}{c|}{$\infty$} & \multicolumn{2}{c|}{$\pi$} \\
\hline
Geodesics & \multicolumn{2}{c|}{Straight lines} & \multicolumn{2}{c|}{Great circles} \\
\hline
Closed geodesics & \multicolumn{2}{c|}{No} & \multicolumn{2}{c|}{Yes (length $2\pi$)} \\
\hline
Parallel lines & \multicolumn{2}{c|}{Never meet} & \multicolumn{2}{c|}{Always meet} \\
\hline
Triangle angle sum & \multicolumn{2}{c|}{$= \pi$} & \multicolumn{2}{c|}{$> \pi$} \\
\hline
Isometry group & $SE(2)$ & $SE(3)$ & $SO(3)$ & $SO(4)$ \\
\hline
Degrees of freedom & 3 & 6 & 3 & 6 \\
\hline
Rotation vs.\ translation & \multicolumn{2}{c|}{Independent} & Identical & Distinct but coupled \\
\hline
Antipodes & \multicolumn{2}{c|}{None} & \multicolumn{2}{c|}{Yes, dist.\ $\pi$} \\
\hline
\end{tabular}
\caption{Comparison of flat and spherical geometries.}
\label{tab:comparison}
\end{table*}

\tableofcontents

\newpage

\section{Introduction}

This document explores the mathematics of flat and spherically curved spaces in two and three dimensions. We develop the geometry of each space, focusing on distances, isometries, and the symmetry groups that describe all possible motions.

Table~\ref{tab:comparison} summarizes the key differences between flat Euclidean space and the spherical spaces $S^2$ and $S^3$.

\section{Flat 2D space}

Let us look at flat two-dimensional space as an example to get familiar with basic geometry and movement through spaces.

We will define flat 2D space as
\[
    A = \left\{ (x, y, 1) \ \mid x, y \in \mathbb{R} \right\} \subset \mathbb{R}^3.
\]

We use a ``light version'' of homogeneous coordinates\footnote{\url{https://en.wikipedia.org/wiki/Homogeneous_coordinates}} with the third dimension fixed to one here. This leads us to the Euclidean group of 2D space without much effort and allows us to follow a top-down approach. It also familiarizes us with homogeneous coordinates, which are heavily used in computer graphics.

We can express transformations of a point $(x, y, 1)$ by multiplying it with an arbitrary $3 \times 3$ matrix from the left:
\[
    \begin{pmatrix} M_{11} & M_{12} & M_{13} \\ M_{21} & M_{22} & M_{23} \\ M_{31} & M_{32} & M_{33}\end{pmatrix}
    \begin{pmatrix}x\\y\\1\end{pmatrix}
    = \begin{pmatrix}M_{11}x + M_{12}y + M_{13}\\M_{21}x + M_{22}y + M_{23}\\M_{31}x + M_{32}y + M_{33}\end{pmatrix}
\]

\subsection{Distance}

The distance between two points $p = (p_x, p_y, 1)$ and $q = (q_x, q_y, 1)$ in flat 2D space is the familiar Euclidean distance:
\[
    d(p, q) = \sqrt{(p_x - q_x)^2 + (p_y - q_y)^2}.
\]

\subsection{Affine group $\mathrm{Aff}(2, \mathbb{R})$}

We are now interested in all matrices that map arbitrary points in $A$ to other points in $A$, or in other words all matrices under which the set $A$ is closed.

This closure requires the last dimension of the resulting vector to be one for any $x$ and $y$:
\[
    M_{31}x + M_{32}y + M_{33} = 1
\]

For this to hold for all $x, y \in \mathbb{R}$, the coefficients of $x$ and $y$ must vanish. The only possible option is $M_{31} = M_{32} = 0$, $M_{33} = 1$.

The affine group $\mathrm{Aff}(2, \mathbb{R})$ consists of all invertible $3 \times 3$ matrices that satisfy the closure requirement introduced above:
\[
    \mathrm{Aff}(2, \mathbb{R}) = \left\{ \begin{pmatrix} M_{11} & M_{12} & M_{13} \\ M_{21} & M_{22} & M_{23} \\ 0 & 0 & 1\end{pmatrix} \;\middle|\; \begin{gathered} M_{ij} \in \mathbb{R} \\ \det M \neq 0 \end{gathered} \right\}
\]

$\mathrm{Aff}(2, \mathbb{R})$ is a subgroup of the general linear group $GL(3, \mathbb{R})$.

\subsection{Special Euclidean group $SE(2)$}

We want to impose two additional restrictions on the set of matrices that we are interested in:
\begin{enumerate}
    \item Isometry: the matrix should preserve distances.
    \item Preservation of handedness: a left hand should remain a left hand, a clockwise angle should remain a clockwise angle, no mirroring allowed.
\end{enumerate}

For a matrix to preserve distances, the upper-left $2 \times 2$ block must be an orthogonal matrix (i.e.\ its columns are orthonormal). Preserving handedness additionally requires $\det R = 1$, which means the block must be a rotation matrix. The general form of an element of $SE(2)$ is therefore:
\[
    SE(2) = \left\{ \begin{pmatrix} \cos\theta & -\sin\theta & t_x \\ \sin\theta & \cos\theta & t_y \\ 0 & 0 & 1 \end{pmatrix} \;\middle|\; \begin{gathered} \theta \in \mathbb{R} \\ t_x, t_y \in \mathbb{R} \end{gathered} \right\}
\]

The group $SE(2)$ is three-dimensional: one degree of freedom for rotation (the angle $\theta$) and two for translation ($t_x, t_y$). Rotation and translation are independent --- any rotation can be combined with any translation.

\subsection{Properties}

Flat 2D space has the familiar geometric properties:
\begin{enumerate}
    \item Straight lines are the shortest paths (geodesics) between two points.
    \item Parallel lines never meet.
    \item The angle sum of a triangle is always $\pi$.
    \item The space extends infinitely in all directions.
\end{enumerate}


\section{Flat 3D space}

Flat, three-dimensional space is what we experience in everyday life. All vectors of three real numbers ($\mathbb{R}^3$) represent distinct points of this space.

\subsection{Distance}

The distance between two points
\[
    p = \begin{pmatrix}p_\mathrm{x}\\p_\mathrm{y}\\p_\mathrm{z}\end{pmatrix} \in \mathbb{R}^3, \quad q = \begin{pmatrix}q_\mathrm{x}\\q_\mathrm{y}\\q_\mathrm{z}\end{pmatrix} \in \mathbb{R}^3
\]
can be calculated as
\[
    d(p, q) = \| q - p \| = \sqrt{(p_\mathrm{x} - q_\mathrm{x})^2 + (p_\mathrm{y} - q_\mathrm{y})^2 + (p_\mathrm{z} - q_\mathrm{z})^2}.
\]

\subsection{Isometries of $\mathbb{R}^3$}

The isometries of $\mathbb{R}^3$ are all transformations $M: \mathbb{R}^3 \to \mathbb{R}^3$ that preserve distances between all points:
\[
    \| M(q) - M(p) \| = \| q - p \|
\]

Permutations of coordinate axes are isometries, for example:
\[
    M(p) = \begin{pmatrix}p_\mathrm{y}\\p_\mathrm{x}\\p_\mathrm{z}\end{pmatrix}
\]

Flipping signs is also okay:
\[
    M(p) = \begin{pmatrix}-p_\mathrm{x}\\p_\mathrm{y}\\-p_\mathrm{z}\end{pmatrix}
\]

Distances are also preserved if we add fixed values to each point. This is called translation. Example:
\[
    M(p) = p + \begin{pmatrix}9\\12\\-50\end{pmatrix}
\]

\subsection{The Euclidean group $E(3)$}

The Euclidean group $E(3)$ contains all isometries of $\mathbb{R}^3$. Therefore, $E(3)$ is the isometry group of $\mathbb{R}^3$. It includes all possible rotations, translations, and mirror reflections in flat 3D space.

The special Euclidean group $SE(3) \subset E(3)$ furthermore excludes mirror reflections (preserves handedness). $SE(3)$ covers all translational and rotational degrees of freedom in 3D space. $SE(3)$ is six-dimensional: three degrees of freedom for translation and three for rotation.

The elements of $SE(3)$ can be described as combinations of rotation matrices (orthogonal matrices with $\det R = 1$) and translation vectors:
\[
    SE(3) = \left\{ (R, t) \;\middle|\; \begin{gathered} R \in \mathbb{R}^{3 \times 3}, R^\top R = I, \det R = 1 \\ t \in \mathbb{R}^3 \end{gathered} \right\}
\]

They act on $\mathbb{R}^3$ with matrix multiplication followed by addition:
\[
    p \mapsto Rp + t
\]

Generally, this action is not a linear transformation but only an affine transformation, since the origin is not preserved if $t$ is not zero.

$SE(3)$ can also be described as a set of 4D matrices that encapsulate rotation matrix and translation vector:
\[
    SE(3) = \left\{ \begin{pmatrix}R_{11}&R_{12}&R_{13}&t_{x}\\R_{21}&R_{22}&R_{23}&t_{y}\\R_{31}&R_{32}&R_{33}&t_{z}\\0&0&0&1\end{pmatrix} \;\middle|\; \begin{gathered} R \in \mathbb{R}^{3 \times 3}, R^\top R = I, \det R = 1 \\ t \in \mathbb{R}^3 \end{gathered} \right\}
\]

With this, the action on $\mathbb{R}^3$ can be reduced to a single matrix multiplication:
\[
    \begin{pmatrix}p_\mathrm{x}\\p_\mathrm{y}\\p_\mathrm{z}\\1\end{pmatrix} \mapsto \begin{pmatrix}R & t \\ 0 & 1\end{pmatrix} \begin{pmatrix}p_\mathrm{x}\\p_\mathrm{y}\\p_\mathrm{z}\\1\end{pmatrix}
\]

We can either let the number 1 in the fourth dimension of our vector appear here out of nowhere, or we can switch to using 4D vectors with a fixed 1 in the fourth dimension. Using such 4D vectors makes it also possible to distinguish between absolute points (with a 1 in the fourth dimension) and relative vectors (with a 0 in the fourth dimension). You can see that the fourth, translational column of the $4 \times 4$ matrix is now only applied to absolute points, preserving the lengths of relative vectors.

More resources on this can be found under the keywords affine space and affine transformations.

Unusual but hopefully demonstrative: We can interpret flat 3D space as being embedded in $\mathbb{R}^4$ with the origin of flat 3D space $(0, 0, 0, 1)$ being the closest point to the 4D origin $(0, 0, 0, 0)$ with a distance of one. (The fourth dimension will be called $w$.)

Interesting, but probably not useful: The points in this embedded flat 3D space with a particular distance $d > 1$ from the 4D origin form a sphere (2-sphere) of radius $\sqrt{d^2 - 1}$ in the embedded 3D space.

\subsection{Quaternion representation of 3D rotations}\label{sec:quaternion-rotation}

Quaternions provide an elegant way to represent 3D rotations without matrices. A quaternion is a number of the form
\[
    q = w + xi + yj + zk
\]
where $i$, $j$, $k$ satisfy the relations $i^2 = j^2 = k^2 = ijk = -1$. The set of all quaternions is denoted $\mathbb{H}$.

A unit quaternion (satisfying $\|q\|^2 = w^2 + x^2 + y^2 + z^2 = 1$) can be written as
\[
    q = \cos\frac{\theta}{2} + \sin\frac{\theta}{2}(u_x i + u_y j + u_z k)
\]
where $\theta$ is a rotation angle and $(u_x, u_y, u_z)$ is the unit axis of rotation.

A rotation can then be applied to a point $p = (p_x, p_y, p_z)$, represented as the pure imaginary quaternion $\hat{p} = p_x i + p_y j + p_z k$, by conjugation:
\[
    \hat{p}' = q \, \hat{p} \, \bar{q}
\]
where $\bar{q} = w - xi - yj - zk$ is the quaternion conjugate. The result $\hat{p}'$ is again a pure imaginary quaternion whose components give the rotated point.

Note that $q$ and $-q$ produce the same rotation, so the unit quaternions form a double cover of $SO(3)$ (i.e., each rotation is represented by exactly two quaternions):
\[
    SO(3) \cong S^3 / \{\pm 1\}
\]

This double cover is the reason the rotation angle appears as $\theta/2$ in the quaternion: a full $2\pi$ rotation of space corresponds to $q$ going from $1$ to $-1$, and a $4\pi$ rotation brings $q$ back to $1$.

Expanding the conjugation $\hat{p}' = q\,\hat{p}\,\bar{q}$ and collecting terms, the rotation corresponding to the unit quaternion $q = w + xi + yj + zk$ can be expressed as the $3 \times 3$ matrix:
\[
    R = \begin{pmatrix}
        1 - 2(y^2 + z^2) & 2(xy - wz)       & 2(xz + wy) \\
        2(xy + wz)       & 1 - 2(x^2 + z^2) & 2(yz - wx) \\
        2(xz - wy)       & 2(yz + wx)        & 1 - 2(x^2 + y^2)
    \end{pmatrix}
\]

In the $4 \times 4$ homogeneous form introduced above, a pure rotation (no translation) becomes:
\[
    \begin{pmatrix}
        1 - 2(y^2 + z^2) & 2(xy - wz)       & 2(xz + wy)       & 0 \\
        2(xy + wz)       & 1 - 2(x^2 + z^2) & 2(yz - wx)        & 0 \\
        2(xz - wy)       & 2(yz + wx)        & 1 - 2(x^2 + y^2) & 0 \\
        0                & 0                 & 0                  & 1
    \end{pmatrix}
\]

This can be verified by checking that $R^\top R = I$ and $\det R = 1$ when $w^2 + x^2 + y^2 + z^2 = 1$.

\subsection{Angles and the dot product}

The angle $\theta$ between two vectors can be found using the dot product:
\[
    \| p \| \cdot \| q \| \cdot \cos \theta = p \cdot q = p_\mathrm{x} \cdot q_\mathrm{x} + p_\mathrm{y} \cdot q_\mathrm{y} + p_\mathrm{z} \cdot q_\mathrm{z}
\]

An important observation for what follows: angles are derived quantities. They are computed from the dot product, which only involves multiplication and addition. When we later work with quaternions and group actions, we can often avoid computing angles explicitly and instead work directly with the algebraic structure. The rotation quaternion $q = \cos\frac{\theta}{2} + \sin\frac{\theta}{2}\hat{u}$ encodes the angle implicitly --- we never need to extract $\theta$ to apply or compose rotations.

\subsection{Properties}

Flat 3D space shares the same familiar properties as its 2D counterpart:
\begin{enumerate}
    \item Geodesics (shortest paths) are straight lines.
    \item Parallel lines never meet.
    \item The space is infinite in all directions.
    \item Rotation and translation are independent: $SE(3)$ decomposes as a semidirect product of $SO(3)$ and $\mathbb{R}^3$.
\end{enumerate}


\section{Spherically curved 2D space}

The 2-sphere $S^2$ is the simplest example of a curved space. It is the two-dimensional surface of a ball in three-dimensional space:
\[
    S^2 = \left\{ (x, y, z) \in \mathbb{R}^3 \mid x^2 + y^2 + z^2 = 1 \right\}
\]

Unlike flat 2D space, $S^2$ is finite (it has surface area $4\pi$) but has no boundary or edge. It is a two-dimensional space that curves through three dimensions.

\subsection{Distance}

The distance between two points $p, q \in S^2$ is the length of the shortest path along the surface, which follows a great circle. Since both points are unit vectors, the geodesic distance is:
\[
    d(p, q) = \arccos(p \cdot q)
\]

The maximum distance between any two points is $\pi$ (half a great circle), achieved by antipodal points $p$ and $-p$.

\subsection{Isometry group: $SO(3)$}

What are the distance-preserving transformations of the 2-sphere? Since $S^2$ is embedded in $\mathbb{R}^3$ as the set of unit vectors, any isometry of $S^2$ must map unit vectors to unit vectors and preserve the dot product (because $d(p,q) = \arccos(p \cdot q)$). These are exactly the orthogonal transformations of $\mathbb{R}^3$ that preserve the sphere.

Restricting to orientation-preserving isometries (excluding reflections), we obtain the rotation group:
\[
    SO(3) = \left\{ R \in \mathbb{R}^{3 \times 3} \mid R^\top R = I, \; \det R = 1 \right\}
\]

$SO(3)$ is three-dimensional: every rotation is determined by an axis (two parameters) and an angle (one parameter).

A key difference from flat 2D space: in flat space, $SE(2)$ has both rotations and translations as independent operations. On the 2-sphere, there is no distinction between rotation and translation. Moving ``forward'' on the sphere is the same as rotating the sphere under your feet. All isometries are rotations. The three degrees of freedom of $SO(3)$ correspond to looking left/right (yaw), looking up/down (pitch), and tilting your head (roll) --- or equivalently, to moving forward/backward, moving left/right, and turning.

\subsection{Stereographic projection}

To visualize the 2-sphere on a flat surface, we can use stereographic projection. Choose a projection point, say the north pole $N = (0, 0, 1)$. For each point $p = (x, y, z) \in S^2$ with $p \neq N$, draw a line from $N$ through $p$ and find where it intersects the $z = 0$ plane:
\[
    (x, y, z) \mapsto \left(\frac{x}{1 - z}, \; \frac{y}{1 - z}\right)
\]

This projection maps the entire sphere (minus the projection point) onto the plane. It has the important property of being conformal: it preserves angles, though not distances or areas. Points near the projection point are mapped far from the origin, and the projection point itself maps to ``infinity.''

Stereographic projection from 4D to 3D, used in this project, works by exact analogy: for a point $(x, y, z, w) \in S^3$, projecting from the point $(0, 0, 0, -1)$:
\[
    (x, y, z, w) \mapsto \left(\frac{x}{w + 1}, \; \frac{y}{w + 1}, \; \frac{z}{w + 1}\right)
\]

\subsection{Properties}

The geometry of $S^2$ differs from flat space in several notable ways:
\begin{enumerate}
    \item Geodesics (shortest paths) are arcs of great circles. Every geodesic is a closed curve of total length $2\pi$.
    \item There are no parallel lines: any two great circles intersect in exactly two (antipodal) points.
    \item The angle sum of a triangle exceeds $\pi$. The excess equals the area of the triangle (Girard's theorem): for a triangle with angles $\alpha$, $\beta$, $\gamma$, the area is $\alpha + \beta + \gamma - \pi$.
    \item Every point has an antipode at distance $\pi$. The antipode of $p$ is $-p$.
\end{enumerate}


\section{Spherically curved 3D space}

The 3-sphere is the three-dimensional analogue of the 2-sphere. Just as $S^2$ is a 2D surface curving through $\mathbb{R}^3$, the 3-sphere $S^3$ is a 3D space curving through $\mathbb{R}^4$.

The 3-sphere can be defined as the set of unit quaternions:
\[
    S^3 = \left\{ a + bi + cj + dk \mid a^2 + b^2 + c^2 + d^2 = 1 \right\}
\]

This identification with unit quaternions is not just a convenient labeling --- the quaternion multiplication gives $S^3$ an algebraic structure (it is a Lie group) that is central to describing its symmetries.

\subsection{Distance}

Since points on $S^3$ are unit quaternions, and $S^3$ is embedded in $\mathbb{R}^4$, the geodesic distance between two points $p, q \in S^3$ is:
\[
    d(p, q) = \arccos(\operatorname{Re}(\bar{p} q))
\]

where $\operatorname{Re}(\bar{p}q)$ is the real (scalar) part of the quaternion product $\bar{p}q$. This is equivalent to the 4D dot product of the two unit vectors. As with $S^2$, the maximum distance is $\pi$, achieved by antipodal points $p$ and $-p$.

In the vertex shader, after applying the view transformation, the $w$-component of the resulting quaternion (in xyzw layout) gives $\operatorname{Re}(\bar{p}q)$ directly, so the spherical distance is computed as $\arccos(w)$.

\subsection{Isometry group: $SO(4)$}

The isometry group of $\mathbb{R}^3$ is the Euclidean group $E(3)$, which can be represented using rotation matrices and translation vectors. However, the isometries of $\mathbb{R}^4$ can move the 3-sphere itself away from the origin. The isometries that are interesting for us are the ones that preserve the 3-sphere. Restricting to orientation-preserving isometries, we obtain $SO(4)$.

The elements of $SO(4)$ perform movement (translation) and rotation of objects that lie on the 3-sphere, such as the camera or models in our demo.

The quaternion structure of $S^3$ gives us a remarkably clean description of $SO(4)$. Every orientation-preserving isometry of $S^3$ can be written as:
\[
    p \mapsto s \cdot p \cdot t
\]
where $s$ and $t$ are unit quaternions. The quaternion $s$ acts by left multiplication and $t$ by right multiplication. This gives a surjective homomorphism
\[
    S^3 \times S^3 \to SO(4), \quad (s, t) \mapsto [p \mapsto s \cdot p \cdot t]
\]

The kernel of this homomorphism is $\{(1, 1), (-1, -1)\}$, since $(-s) \cdot p \cdot (-t) = s \cdot p \cdot t$. This gives the isomorphism:
\[
    SO(4) \cong (S^3 \times S^3) / \mathbb{Z}_2
\]

$SO(4)$ is six-dimensional: the left quaternion $s$ contributes three degrees of freedom, and the right quaternion $t$ contributes another three (each is a unit quaternion on $S^3$, which is three-dimensional).

\subsection{Rotation and translation as special cases}

On the 2-sphere, we noted that rotation and translation are the same thing. On the 3-sphere, the situation is more nuanced: both rotation and translation are isometries of the form $p \mapsto s \cdot p \cdot t$, but they correspond to different relationships between $s$ and $t$.

\textbf{Rotation} (turning in place) fixes the point $1 \in S^3$ (the ``origin''). For $s \cdot 1 \cdot t = 1$, we need $t = \bar{s}$. So a rotation about the origin is:
\[
    p \mapsto s \cdot p \cdot \bar{s}
\]

This is exactly the quaternion rotation formula from Section~\ref{sec:quaternion-rotation}. The three-dimensional family of such rotations forms a copy of $SO(3)$ inside $SO(4)$.

\textbf{Translation} (moving to a new position without turning) is achieved by choosing $t = s$:
\[
    p \mapsto s \cdot p \cdot s
\]

This moves the origin $1$ to $s \cdot 1 \cdot s = s^2$ while preserving the ``forward direction.'' In the demo, pressing the arrow keys applies this type of transformation.

More generally, yaw, pitch, and roll use $(s, t) = (q^{-e}, (-q)^{-e})$ for a basis quaternion $q$ and exponent $e = \alpha/\pi$, while translation uses $(s, t) = (q^{-e}, q^{-e})$.

\subsection{Properties}

The geometry of $S^3$ is a direct generalization of $S^2$ geometry:

\begin{enumerate}
    \item \textbf{Geodesics} are great circles: intersections of $S^3$ with 2-planes through the origin in $\mathbb{R}^4$. Every geodesic is a closed curve of circumference $2\pi$. If you walk in a straight line, you return to your starting point after a distance of $2\pi$.

    \item \textbf{Antipodes.} Every point $p$ has an antipodal point $-p$ at the maximum distance $\pi$. The antipode is the farthest point in all directions --- no matter which direction you walk from $p$, after a distance of $\pi$ you arrive at $-p$.

    \item \textbf{Totally geodesic surfaces} are great 2-spheres: intersections of $S^3$ with 3-dimensional hyperplanes through the origin. These are the analogue of ``planes'' in $S^3$. A great 2-sphere divides $S^3$ into two congruent halves, just as a great circle divides $S^2$ into two hemispheres.

    \item \textbf{Volume.} The total volume of $S^3$ is $2\pi^2$ (with unit radius). This is finite --- the 3-sphere is a closed, bounded space without boundary.

    \item \textbf{Parallel transport.} Unlike flat space, parallel-transporting a vector around a closed loop on $S^3$ can rotate it. This is a manifestation of the curvature of the space.
\end{enumerate}

\end{document}
